\chapter{Mathematical and Financial Background}

%%%%%%%%%%%%%%%%%%%%%%%%%%%%%%%%%%%%%%%%%%%%%%%%%
\section{Mathematical Framework For Option Pricing}
This section defines the common mathematical setup used in the rest of the thesis for pricing European-style derivatives.

Time is continuous on $[0,T]$, with current time $t$ and maturity $T$. Rather than introducing the full probabilistic machinery, we keep a pricing-focused notation that is sufficient for the methods used later in the thesis.

A global notation glossary is provided in Appendix C and will be extended as new symbols are introduced.

As a practical starting point for this thesis, we use the standard risk-neutral valuation formula for European claims with constant rate $r$ \cite[Ch.~8]{oksendal2003}:
\begin{equation}
    V_t=e^{-r\tau}\,\mathbb E_t^{\mathbb Q}[\Psi].
\end{equation}

In the following sections, this framework is connected to three ingredients used throughout the thesis: the stochastic-calculus tools needed for the model equations, the Heston stochastic-volatility dynamics, and the characteristic-function/Fourier-COS pricing route.


%%%%%%%%%%%%%%%%%%%%%%%%%%%%%%%%%%%%%%%%%%%%%%%%%
\section{Stochastic Calculus Background}
This section introduces only the stochastic-calculus notions needed by the pricing methods used later. The presentation follows standard probability definitions and the transform-pricing notation adopted in \cite[Ch.~1]{oksendal2003} and \cite[Sec.~2]{fang_oosterlee_2009}.

For a real random variable $X$, the cumulative distribution function is
\begin{equation}
    F_X(x)=\mathbb P(X\le x).
\end{equation}
If $F_X$ is differentiable, the associated probability density function is
\begin{equation}
    f_X(x)=\frac{\partial F_X(x)}{\partial x}.
\end{equation}

Following this notation, the characteristic function of $X$ is
\begin{equation}
    \phi_X(u) :=\mathbb E\left[ e^{iuX} \right] = \int_{-\infty}^\infty e^{iux}dF_X(x) = \int_{-\infty}^\infty e^{iux}f_X(x)dx,
\end{equation}
\noindent for $u\in\mathbb R$.

Applying the logarithm to $\phi_X$, we obtain the cumulant characteristic function
\begin{equation}
    \zeta_X(u) = \log \mathbb E\left[e^{iux}\right] = \log \phi_X(u).
\end{equation}
Its Maclaurin expansion is
\begin{equation}
    \zeta_X(u) = \sum_{k=0}^\infty \frac{\partial^k \zeta_X (u)}{\partial u^k}\bigg| _{u=0}\frac{u^k}{k!}=\sum_{k=0}^{\infty}\zeta_k \frac{u^k}{k!},
\end{equation}
where $\zeta_k \equiv\frac{\partial^k \zeta_X(u)}{\partial u^k}\bigg|_{u=0}$ is the $k$-th cumulant.


The moment generating function is
\begin{equation}
    \mathcal M_X(u):= \phi_X(-iu) = \mathbb E\left[ e^{ux} \right]. 
\end{equation}


Using $\mathcal M_X(u)=\phi_X(-iu)$ and $\zeta_X(u)=\log \mathcal M_X(u)$, cumulants can be computed from the characteristic function as
\begin{equation}
    \label{eq:cumulants}
    \zeta_k =\frac{\partial ^k \zeta_X(u)}{\partial u^k}\bigg|_{u=0}=\frac{\partial ^k}{\partial u^k}\log \phi_X(-iu)\bigg|_{u=0}.
\end{equation}
These objects are central in the rest of this chapter: the characteristic function provides direct access to Fourier-domain pricing formulas, and low-order cumulants are later used to construct truncation intervals for the COS method.

%%%%%%%%%%%%%%%%%%%%%%%%%%%%%%%%%%%%%%%%%%%%%%%%%
\section{The Heston Stochastic Volatility Model}
In a general diffusion form, the Heston dynamics can be written as
\begin{align}
    dS_t &= \mu S_t\,dt+\sqrt{\nu_t}\,S_t\,dW_t^{(S)},\\
    d\nu_t &= \kappa(\bar\nu-\nu_t)\,dt+\gamma\sqrt{\nu_t}\,dW_t^{(\nu)},\\
    d\langle W^{(S)},W^{(\nu)}\rangle_t &= \rho\,dt.
\end{align}
The parameter $\mu$ denotes a generic drift under the original measure. For pricing, we move to the risk-neutral measure $\mathbb Q$, under which the discounted asset price is a martingale and, in the no-dividend case considered in this thesis, the drift becomes $r$:
\begin{equation}
    dS_t = rS_t\,dt+\sqrt{\nu_t}\,S_t\,dW_t^{(S,\mathbb Q)}.
\end{equation}
This is the form used in the pricing formulas of the following sections \cite{heston,oksendal2003}.

Here, $\nu_t$ is the instantaneous variance (equivalent to the $v_t$ notation used in later data-driven chapters). The variance process is of CIR type, so it is mean-reverting toward $\bar\nu$ at speed $\kappa$, with volatility of variance controlled by $\gamma$.

In this thesis, implied volatility is defined as the value $\sigma_{imp}$ that reproduces a given option price when inserted into the Black--Scholes formula with the same contract and market inputs \cite{black_scholes_1973}:
\begin{equation}
    V_{BS}(S_0,K,r,\tau,\sigma_{imp})=V_{\text{target}}.
\end{equation}
Therefore, implied volatility is a quoting transformation of prices (not a structural parameter of Heston). The numerical inversion used to compute $\sigma_{imp}$ is detailed in Section 2.6.

Parameter interpretation is directly linked to the implied-volatility (IV) surface:
\begin{table}[H]
    \ttabbox[\FBwidth]
    {\caption{Interpretation of Heston parameters}}
    {\begin{tabular}{|P{2.2cm}|P{4.3cm}|P{5.3cm}|}
        \hline
        \textbf{Parameter} & \textbf{Interpretation} & \textbf{Main effect on IV surface} \\
        \hline
        $\rho$ & Instantaneous correlation between spot and variance shocks. & Controls skew: more negative $\rho$ typically implies stronger left skew. \\
        \hline
        $\gamma$ & Volatility of variance (\textit{vol-of-vol}). & Controls curvature and short-maturity smile amplitude; larger $\gamma$ usually increases convexity. \\
        \hline
        $\bar\nu$ & Long-run mean level of the variance process. & Sets the global IV level around which the surface oscillates across maturities. \\
        \hline
        $\kappa$ & Speed of mean reversion of $\nu_t$ toward $\bar\nu$. & Controls how fast maturities converge to long-run variance, affecting term-structure smoothing. \\
        \hline
        \multicolumn{3}{l}{Source: Based on \cite{heston} and \cite[Chs.~2--3]{rouah_heston_book}}
    \end{tabular}}
\end{table}

A common positivity condition for the CIR variance process is the Feller condition
\begin{equation}
    2\kappa\bar\nu \ge \gamma^2.
\end{equation}
In practice, calibrated parameter sets may violate this inequality while still producing usable market fits; therefore, it is treated as a diagnostic and stability indicator rather than as a hard calibration constraint in many numerical workflows \cite[Chs.~2--3]{rouah_heston_book}.

For pricing, the key advantage of Heston is that the log-price characteristic function is available in closed (semi-analytical) form \cite{heston}. This makes the model especially suitable for Fourier-domain valuation, which is the computational route adopted in the next sections.

%%%%%%%%%%%%%%%%%%%%%%%%%%%%%%%%%%%%%%%%%%%%%%%%%
\section{Pricing Via Characteristic Functions}
When the characteristic function is available in closed or semi-closed form, pricing in the Fourier domain becomes a natural alternative to direct density integration. Following \cite[Sec.~2]{fang_oosterlee_2009}, we write the option value in terms of the conditional transition density of the terminal state variable.

Using the notation adopted in the next section, with $x=X(t_0)$ and $y=X(T)$, the discounted valuation formula is
\begin{equation}
    V(t_0,x)=e^{-r\tau}\int_{\mathbb R}V(T,y)\,f_X(T,y;t_0,x)\,dy.
\end{equation}
The corresponding conditional characteristic function is
\begin{equation}
    \phi_X(u;t_0,x):=\mathbb E_{t_0}^{\mathbb Q}\!\left[e^{iuX(T)}\mid X(t_0)=x\right]
    =\int_{\mathbb R}e^{iuy}f_X(T,y;t_0,x)\,dy.
\end{equation}
Under standard integrability conditions, the density can be recovered by Fourier inversion:
\begin{equation}
    f_X(T,y;t_0,x)=\frac{1}{2\pi}\int_{\mathbb R}e^{-iuy}\phi_X(u;t_0,x)\,du.
\end{equation}
Let $\hat V_T(u):=\int_{\mathbb R}e^{-iuy}V(T,y)\,dy$ denote the Fourier transform of the payoff at maturity. Substituting the inverse representation of $f_X$ into the pricing integral and exchanging the order of integration gives
\begin{equation}
    V(t_0,x)=e^{-r\tau}\frac{1}{2\pi}\int_{\mathbb R}\phi_X(u;t_0,x)\,\hat V_T(u)\,du.
\end{equation}

Compared with PDE solvers, this approach avoids full space-time gridding. Compared with Monte Carlo, it is deterministic and does not introduce sampling noise. In repeated-pricing tasks (surface generation, calibration loops, and large synthetic datasets), this is a major practical advantage.

The Fourier-COS method used in the next section is a specific spectral discretization of this transform-based pricing framework: it replaces the continuous integral by a finite cosine expansion over a truncated interval, preserving high accuracy with efficient vectorized evaluation.

%%%%%%%%%%%%%%%%%%%%%%%%%%%%%%%%%%%%%%%%%%%%%%%%%
\section{The Fourier-COS method}

\begin{itemize}
    \item presentar el metodo cos como aproximacion espectral de densidades y precios en un intervalo truncado
    \item comentar con cuidado como elegir el intervalo $[a,b]$ a partir de cumulantes y que impacto tiene esa eleccion en el error (un poco en apendice ya)
    \item separar claramente error de truncacion y error por numero finito de terminos para luego interpretar mejor los resultados numericos (tambien en apendice ahora)
    \item ir diciendo formula, significado de cada termino y como se implementa (vectorizado)
\end{itemize}
\todos{anadir como una guia practica para elegir $N$ y evitar artefactos numericos en extremos de K y vencimiento}

The COS method is the …\\
The main reasons we want to use a cosine expansion are:
\begin{itemize}
    \item In a cosine Fourier expansion we can represent even functions around 0 exactly.
    \item We can extend any function $g:[0,\pi]\to \mathbb R$ to become an even function on $[-\pi, \pi]$ as follows: 
    \begin{equation}
        \bar g(\theta)=\begin{cases}g(\theta), \quad &\theta \geq 0,\\ g(-\theta), \quad &\theta < 0. \end{cases}
    \end{equation}
    \item For functions supported on any other finite interval, say $[a,b]\in\mathbb R$, the Fourier $\cos$ series expansion can be obtained via a change of variables: 
    $$\theta:= \frac{y-a}{b-a}\pi, \quad y:=\frac{b-a}{\pi}\theta+a.$$
\end{itemize}

\todos{TODO: falta más teoría para enlazar}\\

We finally arrive to the approximation of the PDF 
\begin{equation}
    \hat f_X(y) = \sum_{k=0}^{N-1}\!'\  \bar F_k\cos\left(\pi k\frac{y-a}{b-a}\right),
\end{equation}
where $\sum '$ indicates that the first term should be multiplied by $1/2$, and the coefficients depend on the ChF: 
\begin{equation}
    \bar F_k:= \frac{2}{b-a}\operatorname{Re} \left[ \phi_X\left( \frac{\pi k}{b-a} \right) \cdot \exp\left( -\frac{ i\pi ak}{b-a} \right)\right].
\end{equation}

\todos{TODO: no se si poner algo aquí o nueva subsection}\\

The value of a plain vanilla option is given by 
\begin{equation}
    V(t_0,x) = e^{-r\tau}\int_\mathbb R V(T,y)f_X(T,y;t_0,x)dy,
\end{equation}
where $\tau = T-t_0$, $f_X(T,y;t_0,x)$ is the transition probability density of $X(t)$, and $r$ the interest rate.\\

Applying the COS method to the last expression, we arrive to
\begin{equation}
    V(t_0,x)\approx e^{-r\tau} \sum_{k=0}^{N-1}\!'\ \operatorname{Re}\left[ \phi_X\left( \frac{\pi k}{b-a} \right)\cdot \exp \left( -\frac{i\pi ak}{b-a} \right) \right] \cdot H_k,
\end{equation}
with $x$ a function of $S(t_0)$, $\phi_X(u)$ the ChF, and $H_K$ the \textit{payoff coefficients}:
\begin{equation}
    H_k := \frac{2}{b-a}\int_a^b V(T,y) \cos\left( \frac{y-a}{b-a}\pi k \right)dy.
\end{equation}
\todos{TODO: seguir con los payoffs coeffients}

%%%%%%%%%%%%%%%%%%%%%%%%%%%%%%%%%%%%%%%%%%%%%%%%%
\subsection{The COS Method Applied To The Heston Model}
\todos{TODO: rellenar}\\
\begin{itemize}
    \item presentr la teoria anterior al caso heston, con definiciones consistentes de variables y convenciones de signos
    \item verificar que la expresion de $\varphi_H$ que uso coincide con una formulacion estandar y numericamente estable
    \item dejar claro como paso de la formula escalar a evaluacion vectorizada en mallas de strikes y vencimientos
    \item cerrar con una lista corta de sanity checks numericas antes de usar esta formula en benchmarks
\end{itemize}
\todos{revisar simbolos de entrada/salida en esta subseccion para que coincidan exactamente con la implementacion en codigo.}

We need to define the following functions: 
\begin{align}    
    D_1(u)&=\sqrt{\left( \kappa-\gamma \rho iu \right)^2 + \left( u^2+iu \right)\gamma^2},\\
    g(u)&=\frac{\kappa-\gamma \rho iu - D_1}{\kappa - \gamma\rho iu + D_1}.
\end{align}


Using a vector of strikes $\mathbf K$, we find the COS formula for the Heston model:
\begin{equation}
    V(t_0,\mathbf x)\approx \mathbf K e^{-r\tau}\cdot\operatorname{Re} \left\{ \sum_{k=0}^{N-1}\!'\  \varphi_H\left[ \frac{\pi k}{b-a},T; \nu(t_0) \right]U_k\cdot \exp\left[ i\pi k\frac{\mathbf X(t_0)-a}{b-a}\right] \right\},
\end{equation}
where \todos{no se si es la ChF?}
\begin{equation}
\begin{aligned}
    \varphi_H\left(t,T;\nu(t_0)\right) =& \exp\left\{iur\tau + \frac{\nu(t_0)}{\gamma^2}\left(\frac{1-e^{-D_1\tau}}{1-g e^{-D_1\tau}}\right) \left( \kappa-i\rho \gamma u - D_1\right) \right\} \cdot \\
     &\cdot\exp\left\{ \frac{\kappa \bar\nu}{\gamma^2}\left(\tau \left( \kappa-i\rho \gamma u -D_1 \right) -2\log\left( \frac{1-ge^{-D_1\tau}}{1-g} \right) \right)\right\}.    
\end{aligned}
\end{equation}


%%%%%%%%%%%%%%%%%%%%%%%%%%%%%%%%%%%%%%%%%%%%%%%%%
\section{Implied Volatility And Numerical Inversion}
\begin{itemize}
    \item explicar como paso de precios a volatilidad implicita y por que esta inversion es clave para comparar modelos.
    \item documentar el metodo numerico de inversion que uso (por ejemplo, newton o brent), incluyendo tolerancias y criterios de parada.
    \item comentar problemas tipicos de inversion (no convergencia, soluciones no fisicas, sensibilidad a la semilla) y como los trato.
    \item cerrar definiendo claramente que metricas en volatilidad implicita voy a reportar en resultados.
\end{itemize}
\todos{decidir si reporto errores en precio, en iv o ambos segun el objetivo de cada experimento.}

%%%%%%%%%%%%%%%%%%%%%%%%%%%%%%%%%%%%%%%%%%%%%%%%%
\section{State Of The Art}
\begin{itemize}
    \item resumir el estado del arte en tres bloques: metodos clasicos de valoracion, metodos espectrales y modelos sustitutos basados en redes neuronales.
    \item destacar no solo resultados positivos sino tambien limites reportados (generalizacion, estabilidad fuera de distribucion y coste de entrenamiento).
    \item posicionar mi trabajo de forma explicita: que parte reutiliza ideas existentes y que parte aporta una comparacion o implementacion nueva.
    \item cerrar con una tabla comparativa breve que justifique las decisiones metodologicas de los capitulos siguientes.
\end{itemize}
\todos{cerrar el set final de referencias recientes y priorizar articulos con codigo reproducible.}
